% Options for packages loaded elsewhere
\PassOptionsToPackage{unicode}{hyperref}
\PassOptionsToPackage{hyphens}{url}
%
\documentclass[9pt,a4paper]{extarticle}
\usepackage{amsmath,amssymb}
\usepackage[a4paper,left=18mm,right=18mm,top=18mm,bottom=17mm,headheight=14pt]{geometry}
\usepackage{fancyhdr}
\usepackage[spanish,es-nodecimaldot]{babel}
\usepackage{enumitem}
\usepackage{titlesec}
\setlist{nosep,leftmargin=*}
% Distinción visual entre secciones principales
\titleformat{name=\section,numberless}{\normalfont\large\bfseries}{}{0pt}{}
\titlespacing*{\section}{0pt}{16pt plus 2pt minus 1pt}{7pt}
\renewcommand{\refname}{Referencias}
\renewcommand{\arraystretch}{0.98}
\usepackage{iftex}
\ifPDFTeX
  \usepackage[T1]{fontenc}
  \usepackage[utf8]{inputenc}
  \usepackage{textcomp} % provide euro and other symbols
\else % if luatex or xetex
  \usepackage{unicode-math} % this also loads fontspec
  \defaultfontfeatures{Scale=MatchLowercase}
  \defaultfontfeatures[\rmfamily]{Ligatures=TeX,Scale=1}
\fi
\usepackage{lmodern}
\ifPDFTeX\else
  % xetex/luatex font selection
\fi
% Use upquote if available, for straight quotes in verbatim environments
\IfFileExists{upquote.sty}{\usepackage{upquote}}{}
\IfFileExists{microtype.sty}{% use microtype if available
  \usepackage[]{microtype}
  \UseMicrotypeSet[protrusion]{basicmath} % disable protrusion for tt fonts
}{}
\makeatletter
\@ifundefined{KOMAClassName}{% if non-KOMA class
  \IfFileExists{parskip.sty}{%
    \usepackage{parskip}
  }{% else
    \setlength{\parindent}{0pt}
    \setlength{\parskip}{5pt}}
}{% if KOMA class
  \KOMAoptions{parskip=half}}
\makeatother
\usepackage{xcolor}
\usepackage{longtable,booktabs,array}
\setlength{\LTpre}{3pt}
\setlength{\LTpost}{3pt}
\usepackage{calc} % for calculating minipage widths
% Correct order of tables after \paragraph or \subparagraph
\usepackage{etoolbox}
\makeatletter
\patchcmd\longtable{\par}{\if@noskipsec\mbox{}\fi\par}{}{}
\makeatother
% Allow footnotes in longtable head/foot
\IfFileExists{footnotehyper.sty}{\usepackage{footnotehyper}}{\usepackage{footnote}}
\makesavenoteenv{longtable}
\setlength{\emergencystretch}{3em} % prevent overfull lines
\providecommand{\tightlist}{%
  \setlength{\itemsep}{0pt}\setlength{\parskip}{0pt}}
\setcounter{secnumdepth}{-\maxdimen} % remove section numbering
\ifLuaTeX
  \usepackage{selnolig}  % disable illegal ligatures
\fi
\usepackage{bookmark}
\IfFileExists{xurl.sty}{\usepackage{xurl}}{} % add URL line breaks if available
\urlstyle{same}
\hypersetup{
  hidelinks,
  pdfcreator={LaTeX via pandoc}}

\author{}
\date{}

\pagestyle{fancy}
\fancyhf{}
\fancyhead[R]{\footnotesize Entrega 4 (2B) - v1.2}
\fancyfoot[R]{\footnotesize\bfseries \thepage}
\renewcommand{\headrulewidth}{0.4pt}
\renewcommand{\footrulewidth}{0pt}
\begin{document}

% =========================
% CARÁTULA
% =========================
\thispagestyle{empty}
\begin{center}
{\Large\bfseries UNIVERSIDAD TÉCNICA ESTATAL DE QUEVEDO}\\[6pt]
{\normalsize Facultad de Ciencias de la Computación}\\[3pt]
{\normalsize Carrera de Ingeniería de Software}\\[14pt]
{\normalsize Asignatura: Ingeniería de Requerimientos (ISR-401)}\\[18pt]

{\LARGE\bfseries PROTOCOLO DE PRERREGISTRO OSF}\\[8pt]
{\large\bfseries FabroGym - Enfoque 3: Explicabilidad como Requisito No Funcional}\\[5pt]
{\normalsize\itshape Estudio de caso para un componente propuesto de recomendación de rutinas}\\[14pt]
\end{center}

\begin{center}
\begin{tabular}{|>{\bfseries}p{0.25\textwidth}|p{0.66\textwidth}|}
\hline
Proyecto: & FabroGym - Sistema de gestión administrativa y operativa para un gimnasio \\ \hline
Entrega: & Entrega 4 (2B/Defensa Final) - componente empírico \\ \hline
Enfoque seleccionado: & Enfoque 3 - Explicabilidad como Requisito No Funcional \\ \hline
Versión: & 1.2 - ajuste del estado real del estudio previo al registro OSF \\ \hline
Fecha: & 27 de agosto de 2026 \\ \hline
Repositorio: & \url{https://github.com/amorad35/FabroGym_ISR401} \\ \hline
Estado: & Preparado para registro OSF; seis walkthroughs realizados previamente como evidencia formativa \\ \hline
\end{tabular}
\end{center}

\vspace{8pt}
\begin{center}
{\bfseries Equipo PFC - Gimnasio}\\[5pt]
\begin{tabular}{|p{0.44\textwidth}|p{0.47\textwidth}|}
\hline
Alvia Villegas Erick Adalberto & Analista líder / entrevistador \\ \hline
Mera Arias Erick Jhair & Documentador / responsable de encuestas \\ \hline
Mora Duarte Alex José & Modelador / apoyo de análisis \\ \hline
Ponce Rivera Mery Helenmey & Verificadora / calidad de requisitos \\ \hline
Vaca Romero David Octavio & Apoyo documental / evidencias \\ \hline
\end{tabular}
\end{center}

\vfill
\begin{center}
Quevedo, Ecuador - 2026
\end{center}
\clearpage

% =========================
% ÍNDICE
% =========================
\tableofcontents
\clearpage

\begin{longtable}[]{@{}
  >{\raggedright\arraybackslash}p{(\columnwidth - 2\tabcolsep) * \real{0.5000}}
  >{\raggedright\arraybackslash}p{(\columnwidth - 2\tabcolsep) * \real{0.5000}}@{}}
\toprule\noalign{}
\begin{minipage}[b]{\linewidth}\raggedright
\textbf{Campo}
\end{minipage} & \begin{minipage}[b]{\linewidth}\raggedright
\textbf{Detalle}
\end{minipage} \\
\midrule\noalign{}
\endhead
\bottomrule\noalign{}
\endlastfoot
Proyecto & FabroGym - Sistema de gestión de gimnasio \\
Tipo de estudio & Estudio de caso único con componente mixto de validación de requisitos \\
Enfoque de la rúbrica 2B & Enfoque 3 - Explicabilidad como RNF \\
Versión del protocolo & 1.2 \\
Fecha de actualización & 27 de agosto de 2026 \\
Repositorio & github.com/amorad35/FabroGym\_ISR401 \\
Registro previsto & Open Science Framework (OSF) - registro posterior a seis walkthroughs formativos y previo a las actividades analíticas y de validación aún pendientes \\
\end{longtable}

\section*{1. Propósito y relación con la Entrega 4 (2B)}\label{propuxf3sito-y-relaciuxf3n-con-la-entrega-4-2b}
\addcontentsline{toc}{section}{1. Propósito y relación con la Entrega 4 (2B)}

El propósito de este protocolo es documentar y congelar en OSF las decisiones metodológicas y analíticas correspondientes al Enfoque 3 de FabroGym: Explicabilidad como Requisito No Funcional (RNF). El protocolo se alinea con la adaptación temática de la rúbrica de la Entrega 4 (2B), que asigna a FabroGym el Enfoque 3 (explicabilidad) \cite{uteq2026}.

Antes del registro OSF ya se realizaron seis sesiones walkthrough: tres con participantes de perfil técnico y tres con participantes de perfil no técnico. Estas sesiones constituyen evidencia empírica previa o formativa y se utilizarán de manera transparente para elicitar, identificar y refinar necesidades y candidatos a RNF de explicabilidad. No se presentarán como datos confirmatorios preregistrados ni se afirmará que ocurrieron después del sello temporal de OSF.

El registro OSF se utilizará para dejar fijadas las decisiones de análisis y las actividades que todavía se encuentren pendientes, entre ellas la sistematización final de la evidencia, la validación de los candidatos a RNF y, si aún no se ha ejecutado, una sesión de member checking con participantes previos del estudio.

El estudio no presupone que el componente de inteligencia artificial o aprendizaje automático ya esté implementado. Su objeto es determinar qué explicaciones deberían exigirse como RNF verificables y cómo son percibidas por perfiles técnicos y no técnicos antes de consolidar el alcance terminal del ERS/SRS y del MVP.

\section*{2. Preguntas de investigación (RQ) y PICOC}\label{preguntas-de-investigaciuxf3n-rq-y-picoc}
\addcontentsline{toc}{section}{2. Preguntas de investigación (RQ) y PICOC}

\begin{longtable}[]{@{}
  >{\raggedright\arraybackslash}p{(\columnwidth - 2\tabcolsep) * \real{0.5000}}
  >{\raggedright\arraybackslash}p{(\columnwidth - 2\tabcolsep) * \real{0.5000}}@{}}
\toprule\noalign{}
\begin{minipage}[b]{\linewidth}\raggedright
\textbf{ID}
\end{minipage} & \begin{minipage}[b]{\linewidth}\raggedright
\textbf{Pregunta}
\end{minipage} \\
\midrule\noalign{}
\endhead
\bottomrule\noalign{}
\endlastfoot
RQ1 & ¿Qué requisitos de explicabilidad emergen para los perfiles técnicos y no técnicos de FabroGym ante un componente de recomendación de rutinas? \\
RQ2 & ¿Existen diferencias en la percepción de la explicabilidad entre usuarios técnicos y no técnicos cuando evalúan las mismas explicaciones del componente de recomendación? \\
\end{longtable}

\begin{longtable}[]{@{}
  >{\raggedright\arraybackslash}p{(\columnwidth - 2\tabcolsep) * \real{0.5000}}
  >{\raggedright\arraybackslash}p{(\columnwidth - 2\tabcolsep) * \real{0.5000}}@{}}
\toprule\noalign{}
\begin{minipage}[b]{\linewidth}\raggedright
\textbf{Elemento PICOC}
\end{minipage} & \begin{minipage}[b]{\linewidth}\raggedright
\textbf{Operacionalización}
\end{minipage} \\
\midrule\noalign{}
\endhead
\bottomrule\noalign{}
\endlastfoot
P - Población & Personas vinculadas al dominio FabroGym: perfiles técnicos de software/IR y perfiles no técnicos del gimnasio o usuarios del servicio, reclutados con consentimiento informado. \\
I - Intervención & Exposición a escenarios y prototipos/mockups que muestran una recomendación de rutina junto con una explicación estructurada de por qué fue sugerida. \\
C - Comparación & Contraste descriptivo y cualitativo entre las necesidades de explicabilidad identificadas en participantes técnicos y no técnicos. \\
O - Resultados & Necesidades y candidatos a RNF de explicabilidad derivados de los walkthroughs; cobertura de dimensiones cuando pueda determinarse con la evidencia disponible; diferencias descriptivas entre perfiles; y resultado del member checking si se ejecuta. \\
C - Contexto & FabroGym, sistema de gestión de gimnasio local en Ecuador, con componente propuesto de recomendación de rutinas. \\
\end{longtable}

\section*{3. Proposiciones de análisis}\label{proposiciones-de-analisis}
\addcontentsline{toc}{section}{3. Proposiciones de análisis}

Para RQ1 y RQ2 se utilizarán proposiciones de estudio de caso, coherentes con la naturaleza cualitativa y descriptiva de las seis sesiones walkthrough disponibles:

\begin{itemize}
\item
  P1: Los perfiles técnicos y no técnicos pueden priorizar necesidades distintas de explicación; estas diferencias se examinarán mediante contraste descriptivo y cualitativo de la evidencia disponible.
\item
  P2: Una explicación útil para FabroGym deberá traducirse en requisitos verificables, no solo en preferencias subjetivas de interfaz.
\item
  P3: Los candidatos a RNF derivados de los walkthroughs deberán conservar trazabilidad con la evidencia que sustenta su formulación y podrán ser confirmados, ajustados o descartados durante el member checking.
\end{itemize}

No se preregistran hipótesis inferenciales para los seis walkthroughs existentes, debido a que la evidencia disponible corresponde a tres participantes técnicos y tres no técnicos y no se han declarado respuestas Likert adicionales como parte de esta base empírica.

\section*{4. Constructos, variables y operacionalización}\label{constructos-variables-y-operacionalizaciuxf3n}
\addcontentsline{toc}{section}{4. Constructos, variables y operacionalización}

\begin{longtable}[]{@{}
  >{\raggedright\arraybackslash}p{(\columnwidth - 8\tabcolsep) * \real{0.2000}}
  >{\raggedright\arraybackslash}p{(\columnwidth - 8\tabcolsep) * \real{0.2000}}
  >{\raggedright\arraybackslash}p{(\columnwidth - 8\tabcolsep) * \real{0.2000}}
  >{\raggedright\arraybackslash}p{(\columnwidth - 8\tabcolsep) * \real{0.2000}}
  >{\raggedright\arraybackslash}p{(\columnwidth - 8\tabcolsep) * \real{0.2000}}@{}}
\toprule\noalign{}
\begin{minipage}[b]{\linewidth}\raggedright
\textbf{Variable / constructo}
\end{minipage} & \begin{minipage}[b]{\linewidth}\raggedright
\textbf{Rol}
\end{minipage} & \begin{minipage}[b]{\linewidth}\raggedright
\textbf{Operacionalización}
\end{minipage} & \begin{minipage}[b]{\linewidth}\raggedright
\textbf{Fuente}
\end{minipage} & \begin{minipage}[b]{\linewidth}\raggedright
\textbf{Escala / tratamiento}
\end{minipage} \\
\midrule\noalign{}
\endhead
\bottomrule\noalign{}
\endlastfoot
Perfil de participante & Comparación & Técnico / no técnico, conforme a la caracterización documentada de cada sesión & Ficha o evidencia de caracterización disponible & Nominal \\
Necesidad de explicabilidad & Analítica & Unidad temática expresada por el participante sobre qué debe explicar el sistema, por qué y con qué nivel de detalle & Registro o transcripción del walkthrough & Categórica / cualitativa \\
Candidato a RNF & Resultado de ingeniería & Enunciado verificable derivado de una o más necesidades de explicabilidad con trazabilidad a su fuente & Matriz de candidatos a RNF & Cualitativa / verificable \\
Cobertura de explicabilidad & Descriptiva & Dimensiones de explicabilidad representadas por al menos un candidato a RNF cuando la codificación disponible permita establecerlo & Matriz RNF-dimensión & Conteo y proporción descriptiva \\
Decisión de member checking & Validación & Confirmado, ajustado o no confirmado para cada interpretación o candidato revisado, si se ejecuta la sesión pendiente & Acta de member checking & Nominal descriptiva \\
\end{longtable}

Las dimensiones empleadas para organizar la evidencia deberán derivarse de la literatura de explicabilidad citada por la rúbrica y de la codificación trazable de los walkthroughs \cite{chazette2020,chazette2021,chazette2022}. No se incorporarán puntuaciones Likert, coeficientes de acuerdo ni variables cuantitativas que no correspondan a datos realmente disponibles.

\section*{5. Diseño del estudio}\label{diseuxf1o-del-estudio}
\addcontentsline{toc}{section}{5. Diseño del estudio}

Se emplea un estudio de caso único con unidades de análisis embebidas por perfil de participante, siguiendo lineamientos metodológicos de estudios de caso en ingeniería de software \cite{runeson2009,wohlin2012}. La evidencia empírica disponible está constituida por seis walkthroughs reales: tres realizados con participantes técnicos y tres con participantes no técnicos.

Los seis walkthroughs fueron ejecutados antes del sello temporal de OSF. Por tanto, se consideran evidencia empírica previa o formativa y no datos confirmatorios preregistrados. El registro OSF no se utilizará para retrodatarlos, sino para documentar su existencia y congelar las decisiones de análisis y validación que aún estén pendientes.

La secuencia metodológica correspondiente al estado real del estudio es la siguiente:

\begin{enumerate}
\def\labelenumi{\arabic{enumi}.}
\item
  Conservar y organizar la evidencia de los seis walkthroughs ya realizados, manteniendo la distinción entre tres participantes técnicos y tres no técnicos.
\item
  Sistematizar las observaciones disponibles para identificar necesidades de explicabilidad y formular o refinar candidatos a RNF verificables con trazabilidad a su evidencia de origen.
\item
  Registrar en OSF el protocolo actualizado, declarando expresamente que los walkthroughs preceden al sello temporal del registro.
\item
  Ejecutar únicamente las actividades que todavía estén pendientes y que hayan quedado definidas en este protocolo, sin declarar nuevas rondas de walkthrough inexistentes.
\item
  Realizar, si todavía no se ha ejecutado, una sesión final de member checking con al menos tres participantes previos del estudio para contrastar la interpretación de los hallazgos y candidatos a RNF.
\item
  Analizar de manera reproducible la evidencia realmente disponible y reportar como no aplicable o no calculable cualquier análisis que requiera datos inexistentes.
\end{enumerate}

El diseño no incorpora rondas adicionales de walkthrough posteriores al OSF ni divide las seis sesiones existentes en rondas artificiales.

\section*{6. Participantes, muestreo y criterios}\label{participantes-muestreo-y-criterios}
\addcontentsline{toc}{section}{6. Participantes, muestreo y criterios}

\begin{longtable}[]{@{}
  >{\raggedright\arraybackslash}p{(\columnwidth - 2\tabcolsep) * \real{0.5000}}
  >{\raggedright\arraybackslash}p{(\columnwidth - 2\tabcolsep) * \real{0.5000}}@{}}
\toprule\noalign{}
\begin{minipage}[b]{\linewidth}\raggedright
\textbf{Aspecto}
\end{minipage} & \begin{minipage}[b]{\linewidth}\raggedright
\textbf{Situación documentada / tratamiento}
\end{minipage} \\
\midrule\noalign{}
\endhead
\bottomrule\noalign{}
\endlastfoot
Perfiles & Técnico: experiencia vinculada a software, ingeniería de requisitos, QA, desarrollo o docencia técnica. No técnico: participante del dominio del gimnasio sin rol profesional de software en el estudio. \\
Walkthroughs realizados & 6 sesiones: 3 con participantes técnicos y 3 con participantes no técnicos. Estas sesiones fueron ejecutadas antes del registro OSF. \\
Member checking & Actividad final prevista por la rúbrica con al menos 3 participantes previos del estudio, únicamente si todavía no ha sido ejecutada. No requiere reclutar necesariamente participantes nuevos. \\
Unidad de análisis & Observación, comentario o decisión atribuible a un participante durante el walkthrough y su relación trazable con una necesidad o candidato a RNF de explicabilidad. \\
\end{longtable}

La elegibilidad de cada sesión para el análisis se verificará con la evidencia documental disponible. No se asumirá ni declarará consentimiento, caracterización, fecha u otro dato que no pueda verificarse. Cuando algún elemento requerido no pueda comprobarse, se registrará como no verificable y se informará de forma transparente en el análisis.

\section*{7. Instrumentos y materiales}\label{instrumentos-y-materiales}
\addcontentsline{toc}{section}{7. Instrumentos y materiales}

\begin{itemize}
\item
  Guion o guía de walkthrough disponible para documentar el procedimiento y las observaciones de explicabilidad; su correspondencia con cada sesión se verificará con la evidencia existente.
\item
  Prototipo o mockup del componente de recomendación de rutinas presentado durante las sesiones, identificado con la versión que pueda verificarse en la evidencia disponible.
\item
  Evidencia disponible de los seis walkthroughs realizados, incluidas notas, registros o transcripciones anonimizadas cuando existan.
\item
  Matriz de candidatos a RNF de explicabilidad: ID, enunciado, fuente, perfil, dimensión, criterio de aceptación y observaciones, completada únicamente con información trazable a la evidencia disponible.
\item
  Ficha de caracterización del participante, cuando exista o pueda reconstruirse legítimamente a partir de información ya documentada sin incorporar identificadores innecesarios.
\item
  Acta de member checking para la actividad final, si esta aún se encuentra pendiente.
\item
  Scripts de importación, validación y análisis reproducible ajustados a la estructura real de los datos disponibles.
\end{itemize}

La estructura de los instrumentos refleja la secuencia real del estudio y no emplea campos de ronda. Los instrumentos no contendrán resultados, puntuaciones o decisiones que no hayan sido obtenidos empíricamente.

\section*{8. Procedimiento de recolección y actividades pendientes}\label{procedimiento-de-recolecciuxf3n}
\addcontentsline{toc}{section}{8. Procedimiento de recolección y actividades pendientes}

\textbf{Actividades realizadas.} Se ejecutaron seis sesiones walkthrough antes del registro OSF: tres con participantes técnicos y tres con participantes no técnicos. Durante la sistematización se conservará la fecha real de ejecución que conste en cada evidencia y no se modificará su clasificación temporal por la fecha de carga a GitHub u OSF.

\textbf{Tratamiento posterior al walkthrough.} La evidencia disponible se anonimizará o seudonimizará según corresponda y se codificará para identificar necesidades de explicabilidad y candidatos a RNF. Cada candidato deberá mantener trazabilidad con la observación o evidencia que lo sustenta.

\textbf{Registro OSF.} El protocolo actualizado se registrará después de los seis walkthroughs. El registro documentará de forma explícita este hecho y congelará las decisiones analíticas y actividades pendientes, sin presentar las sesiones anteriores como preregistradas.

\textbf{Member checking.} Si esta actividad aún no se ha ejecutado, se realizará una sesión final con al menos tres participantes previos del estudio. La sesión se utilizará para confirmar, ajustar o no confirmar las interpretaciones y candidatos a RNF derivados de la evidencia.

\textbf{Cierre y análisis.} Después de completar las actividades pendientes se ejecutará el análisis sobre la evidencia real disponible. Cualquier análisis para el cual no existan datos suficientes se declarará como no aplicable o no calculable.

\section*{9. Plan de análisis preregistrado}\label{plan-de-anuxe1lisis-preregistrado}
\addcontentsline{toc}{section}{9. Plan de análisis preregistrado}

El análisis se limitará a la evidencia realmente disponible y se ejecutará mediante procedimientos reproducibles. No se incorporarán métricas inferenciales, respuestas Likert, intervalos de confianza ni coeficientes de acuerdo si los datos necesarios no existen o no pueden verificarse.

\begin{itemize}
\item
  RQ1: codificación y síntesis cualitativa de las necesidades de explicabilidad expresadas en los seis walkthroughs; formulación y refinamiento de candidatos a RNF con trazabilidad a la evidencia de origen.
\item
  RQ2: contraste descriptivo y cualitativo entre los patrones identificados en los tres participantes técnicos y los tres no técnicos. Se podrán reportar conteos de categorías o dimensiones únicamente cuando deriven de una codificación reproducible de la evidencia.
\item
  Cobertura de explicabilidad: número y proporción de dimensiones representadas por al menos un candidato a RNF verificable, únicamente cuando la matriz de codificación permita establecer de forma trazable cuáles dimensiones fueron evaluadas.
\item
  Member checking: resumen descriptivo de candidatos o interpretaciones confirmados, ajustados o no confirmados durante la sesión final, si esta se ejecuta.
\item
  No se prevé calcular kappa entre rondas, al no existir rondas reales comparables en el estudio. Un coeficiente de acuerdo solo podrá considerarse si posteriormente existe una matriz nominal real con evaluadores comparables y esta situación se documenta como análisis adicional.
\item
  No se aplicará U de Mann-Whitney ni se calcularán tamaños del efecto o intervalos de confianza sobre puntuaciones Likert mientras no existan respuestas Likert reales y verificables correspondientes al estudio de explicabilidad.
\item
  Cualquier análisis no previsto o no sustentado por la evidencia disponible se etiquetará como exploratorio, no aplicable o no calculable, según corresponda.
\end{itemize}

\clearpage
\section*{10. Gestión de datos, ética y reproducibilidad}\label{gestiuxf3n-de-datos-uxe9tica-y-reproducibilidad}
\addcontentsline{toc}{section}{10. Gestión de datos, ética y reproducibilidad}

\begin{itemize}
\item
  Zona pública {[}P{]}: datos anonimizados, transcripciones seudonimizadas, instrumentos vacíos, matrices procesadas y scripts.
\item
  Zona restringida {[}R{]}: consentimientos originales, firmas, cédulas, audios, videos y cualquier otro identificador directo; almacenamiento cifrado AES-256 conforme a la rúbrica.
\item
  No se publicarán datos de salud, biométricos ni identificadores reales de clientes.
\item
  Los resultados del manuscrito deberán poder regenerarse desde los datos crudos anonimizados mediante los scripts versionados.
\item
  El registro OSF conservará el protocolo original; las modificaciones posteriores se documentarán en osf\_deviations.pdf y en la sección de desviaciones del registro, en coherencia con los principios de prerregistro \cite{nosek2018}.
\end{itemize}

\section*{11. Tratamiento de actividades realizadas antes del registro OSF}\label{tratamiento-de-actividades-realizadas-antes-del-registro-osf}
\addcontentsline{toc}{section}{11. Tratamiento de actividades realizadas antes del registro OSF}

Antes del registro OSF se realizaron seis walkthroughs reales: tres con participantes técnicos y tres con participantes no técnicos. Estas seis sesiones constituyen evidencia empírica previa o formativa del estudio y satisfacen el mínimo de sesiones de walkthrough requerido por la rúbrica terminal 2B para este componente de validación.

Debido a que las sesiones preceden al sello temporal del OSF, no se describirán como datos confirmatorios preregistrados. El registro OSF no altera su fecha ni su naturaleza metodológica. Su utilización se limita a la elicitación, identificación y refinamiento de necesidades y candidatos a RNF de explicabilidad, así como al contraste descriptivo entre perfiles.

La fecha de incorporación de archivos a GitHub u OSF no modifica la fecha real de ejecución de las sesiones. La trazabilidad deberá conservar la secuencia temporal efectiva del estudio.

El registro OSF congelará las decisiones analíticas y las actividades aún pendientes. Entre estas podrá incluirse el member checking con al menos tres participantes previos, si todavía no ha sido realizado. No se declararán nuevas rondas de walkthrough ni se atribuirán al prerregistro actividades ya ejecutadas.

\section*{12. Reglas de datos faltantes, evidencia no verificable y exclusiones}\label{reglas-de-datos-faltantes-y-exclusiones}
\addcontentsline{toc}{section}{12. Reglas de datos faltantes, evidencia no verificable y exclusiones}

\begin{itemize}
\item
  No se imputarán respuestas, observaciones, fechas, perfiles, consentimientos ni decisiones que no consten en la evidencia disponible.
\item
  Cuando una sesión carezca de información necesaria para un análisis específico, esa ausencia se registrará como dato no disponible o no verificable y se reportará el número efectivo de casos utilizados.
\item
  Una observación solo se atribuirá a un perfil técnico o no técnico cuando la clasificación pueda verificarse en la documentación disponible.
\item
  Los registros duplicados verificables se depurarán manteniendo trazabilidad de la decisión adoptada.
\item
  Toda exclusión o imposibilidad de cálculo deberá quedar documentada de forma reproducible.
\end{itemize}

\section*{13. Desviaciones y cambios al protocolo}\label{desviaciones-y-cambios-al-protocolo}
\addcontentsline{toc}{section}{13. Desviaciones y cambios al protocolo}

Cualquier cambio posterior al registro se documentará con: (a) descripción del cambio, (b) razón, (c) fecha en que se detectó la necesidad, (d) si ocurrió antes o después de observar resultados, y (e) mitigación. Los análisis no preregistrados se identificarán como exploratorios. La existencia de una desviación no implica ocultarla ni reescribir el protocolo original.

\section*{14. Productos esperados}\label{productos-esperados}
\addcontentsline{toc}{section}{14. Productos esperados}

\begin{longtable}[]{@{}
  >{\raggedright\arraybackslash}p{(\columnwidth - 2\tabcolsep) * \real{0.5000}}
  >{\raggedright\arraybackslash}p{(\columnwidth - 2\tabcolsep) * \real{0.5000}}@{}}
\toprule\noalign{}
\begin{minipage}[b]{\linewidth}\raggedright
\textbf{Archivo / ubicación}
\end{minipage} & \begin{minipage}[b]{\linewidth}\raggedright
\textbf{Contenido}
\end{minipage} \\
\midrule\noalign{}
\endhead
\bottomrule\noalign{}
\endlastfoot
06\_Experimento/protocolo.pdf & Versión del protocolo correspondiente al estado real del estudio al momento del registro OSF. \\
06\_Experimento/osf\_registration.pdf & Comprobante oficial del registro OSF con URL persistente y sello temporal. \\
06\_Experimento/osf\_deviations.pdf & Desviaciones documentadas respecto del protocolo, si las hubiera. \\
06\_Experimento/instrumentos/ & Instrumentos y matrices utilizados para sistematizar los walkthroughs y realizar member checking cuando corresponda. \\
06\_Experimento/resultados/ & Tablas, matrices y figuras derivadas de la evidencia real disponible mediante procedimientos reproducibles. \\
06\_Experimento/scripts\_analisis/ & Scripts ajustados a la estructura real de los datos utilizados en el análisis final. \\
\end{longtable}

\section*{15. Resumen listo para copiar en OSF}\label{resumen-listo-para-copiar-en-osf}
\addcontentsline{toc}{section}{15. Resumen listo para copiar en OSF}

\begin{longtable}[]{@{}
  >{\raggedright\arraybackslash}p{(\columnwidth - 2\tabcolsep) * \real{0.5000}}
  >{\raggedright\arraybackslash}p{(\columnwidth - 2\tabcolsep) * \real{0.5000}}@{}}
\toprule\noalign{}
\begin{minipage}[b]{\linewidth}\raggedright
\textbf{Campo OSF}
\end{minipage} & \begin{minipage}[b]{\linewidth}\raggedright
\textbf{Texto recomendado}
\end{minipage} \\
\midrule\noalign{}
\endhead
\bottomrule\noalign{}
\endlastfoot
Título & FabroGym: registration of an explainability requirements analysis for a fitness routine recommendation component \\
Descripción & Estudio de caso sobre requisitos no funcionales de explicabilidad en un componente propuesto de recomendación de rutinas. Antes del registro OSF se realizaron seis walkthroughs: tres con participantes técnicos y tres con participantes no técnicos. Estas sesiones se tratan como evidencia empírica previa/formativa para elicitar y refinar necesidades y candidatos a RNF. El registro congela las decisiones de análisis y las actividades de validación todavía pendientes. \\
Pregunta principal & ¿Qué requisitos de explicabilidad emergen para los perfiles técnicos y no técnicos de FabroGym y qué diferencias descriptivas se observan entre ambos perfiles? \\
Diseño & Estudio de caso único basado en seis walkthroughs previos al OSF, con análisis cualitativo y descriptivo de la evidencia y member checking final si aún se encuentra pendiente. \\
Análisis principal & Codificación cualitativa de necesidades de explicabilidad; formulación y trazabilidad de candidatos a RNF; contraste descriptivo entre perfiles; cobertura de dimensiones cuando sea calculable; y síntesis del member checking si se ejecuta. \\
Estado al registrar & Seis walkthroughs ya realizados antes del sello temporal de OSF: 3 técnicos y 3 no técnicos. No se presentan como datos confirmatorios preregistrados. No se declaran rondas adicionales inexistentes ni análisis que requieran datos no disponibles. \\
\end{longtable}

\addcontentsline{toc}{section}{Referencias}

\begin{thebibliography}{99}

\bibitem{chazette2020}
L. Chazette and K. Schneider, ``Explainability as a non-functional requirement: Challenges and recommendations,'' \textit{Requirements Engineering}, vol. 25, pp. 493--514, 2020.

\bibitem{chazette2021}
L. Chazette, W. Brunotte, and T. Speith, ``Exploring explainability: A definition, a model, and a knowledge catalogue,'' in \textit{Proc. IEEE Int. Requirements Engineering Conf. (RE)}, 2021, pp. 197--208.

\bibitem{chazette2022}
L. Chazette, W. Brunotte, and T. Speith, ``Explainable software systems: From requirements analysis to system evaluation,'' \textit{Requirements Engineering}, vol. 27, no. 4, pp. 457--487, 2022.

\bibitem{nosek2018}
B. A. Nosek, C. R. Ebersole, A. C. DeHaven, and D. T. Mellor, ``The preregistration revolution,'' \textit{Proc. Natl. Acad. Sci. U.S.A.}, vol. 115, no. 11, pp. 2600--2606, 2018, doi: 10.1073/pnas.1708274114.

\bibitem{runeson2009}
P. Runeson and M. Höst, ``Guidelines for conducting and reporting case study research in software engineering,'' \textit{Empirical Software Engineering}, vol. 14, no. 2, pp. 131--164, 2009, doi: 10.1007/s10664-008-9102-8.

\bibitem{wohlin2012}
C. Wohlin, P. Runeson, M. Höst, M. C. Ohlsson, B. Regnell, and A. Wesslén, \textit{Experimentation in Software Engineering}. Berlin, Germany: Springer, 2012.

\bibitem{uteq2026}
Universidad Técnica Estatal de Quevedo, ``Guía y rúbrica de evaluación - Proyecto Fin de Curso, Entrega 4 (2B/Defensa Final): Manuscrito para revista JCR,'' Ingeniería de Requerimientos ISR-401, 2026.

\end{thebibliography}

\end{document}
